\SourceLeaf{118}
\Article{24}
Die vorangehenden Versuche wurden auf der Sternwarte durchgeführt. Die Apparate wurden so aufgestellt, dass Eisen möglichst weit von ihnen ferngehalten wurde. Trotzdem ist nicht zu bezweifeln, dass die zahlreichen Eisenmassen in den Wänden, Fenstern und Türen des Gebäudes und selbst die Eisenteile der größeren astronomischen Instrumente, in denen die erdmagnetische Kraft selbst Magnetismus hervorruft, eine merkliche Wirkung auf die aufgehängten Nadeln ausüben. Die daraus entstehenden Kräfte verändern sowohl die Richtung als auch die Intensität des Erdmagnetismus ein wenig. Unsere Versuche liefern daher den am Standort des Apparats $A$ durch diese Störung veränderten Wert der Intensität. Diese Abweichung muss praktisch konstant bleiben, solange die Eisenmassen ihre Lage beibehalten und sich die Elemente des Erdmagnetismus selbst, nämlich Intensität und Richtung, nicht wesentlich ändern. Wie groß die Veränderung ist, ist bisher unbekannt, doch würde ich kaum glauben, dass sie über ein oder zwei Hundertstel des Gesamtwerts hinausgeht. Es dürfte jedoch nicht schwierig sein, ihre Größe wenigstens näherungsweise durch Versuche zu bestimmen. Dazu wären die gleichzeitigen Schwingungen zweier Nadeln zu beobachten, von denen die eine am gewohnten Ort auf der Sternwarte, die andere im Freien in hinreichender Entfernung vom Gebäude oder von anderen störenden Eisenmassen aufgehängt würde. Anschließend müssten die Nadeln ihre Plätze tauschen. Bisher hat sich keine Gelegenheit ergeben, diese Versuche durchzuführen. Die sicherste Abhilfe wird jedoch ein eigenes Gebäude für magnetische Beobachtungen schaffen, das dank königlicher Großzügigkeit bald errichtet werden soll. Bei seinem Bau wird Eisen vollständig ausgeschlossen.

